% !TEX program = xelatex
\documentclass[hyperref,a4paper,UTF8]{ctexart}

\usepackage[left=2.50cm, right=2.50cm, top=2.50cm, bottom=2.50cm]{geometry}

\usepackage[unicode=true,colorlinks,urlcolor=blue,linkcolor=blue,bookmarksnumbered=true]{hyperref}
\usepackage{latexsym,amssymb,amsmath,amsbsy,amsopn,amstext,amsthm,amsxtra,color,bm,calc,ifpdf}
\usepackage{graphicx}
\usepackage{enumerate}
\usepackage{fancyhdr}
\usepackage{listings}
\usepackage{multirow}
\usepackage{makeidx}
\usepackage{xcolor}
\usepackage{fontspec}
\usepackage{subfigure}
\usepackage{hyperref}
\usepackage{pythonhighlight}

\pagestyle{fancy}
\fancyhead[L]{}
\fancyhead[C]{\fangsong 在此填写页眉（可以是论文题目）}
\fancyhead[R]{}

\renewcommand{\abstractname}{\textbf{\large {摘\quad 要}}} % 更改摘要二字的样式


\title{\textbf{{在此填写题目}}}
\author{
\kaishu\normalsize
姓名\ \underline{在此填写姓名} \qquad
学号\ \underline{在此填写学号} \qquad
院系\ \underline{在此填写院系}
}
\date{} % 留空，不显示日期


\begin{document}

\begin{figure}
    \centering
    \includegraphics[width=0.35\textwidth]{figures/SISU.png}
\end{figure}

\maketitle

\begin{abstract}

在此填写摘要内容

\end{abstract}

\
\newpage
\tableofcontents

\thispagestyle{empty} % 当前页不显示页码
\newpage


\section{引言}

在此填写引言

\

\section{第2章标题}

以下是一些样例：\cite{Author2020}


\section{结论}

在此填写结论

\


\bibliography{references} % 使用你的.bib文件名，不包含扩展名
\bibliographystyle{plain} % 选择引用样式，可以根据需求选择不同的样式


\end{document}